\documentclass[11pt]{article}
\usepackage[margin=1in]{geometry}
\usepackage{amsmath,amssymb,amsthm,mathtools}
\usepackage{graphicx}
\usepackage{microtype}
\usepackage[hidelinks]{hyperref}

\newcommand{\R}{\mathbb R}
\newcommand{\N}{\mathbb N}
\newcommand{\G}{\Gamma}
\newcommand{\bi}{\mathbin{\boxtimes_{i}}}
\newtheorem{theorem}{Theorem}
\newtheorem{lemma}{Lemma}

\title{An integral tensor outside the algebraic completion}
\author{Candidate counterexample for arXiv:2404.16338}
\date{19 July 2026}

\begin{document}
\maketitle

\begin{center}
\fbox{\parbox{0.91\textwidth}{
\textbf{Status: candidate counterexample, likely valid.}
The integral-representation space of Definition 4.1 need not be the
completion of the algebraic tensor product for its defining seminorms.  The
counterexample is explicit and uses two copies of a step-function model of
$\ell_\infty$.}}
\end{center}

\section{Source and question}

The source is Eva-Maria Hekkelman, Edward McDonald, and Teun D. H. van
Nuland, \emph{Multiple operator integrals, pseudodifferential calculus, and
asymptotic expansions}, arXiv:2404.16338, 2024 \cite{HMN}.  In Definition 4.1
they introduce $\G_0\bi\cdots\bi\G_n$: its elements are functions admitting
representations
\begin{equation}\label{eq:source-rep}
 \phi(\lambda_0,\ldots,\lambda_n)
 =\int_\Omega \prod_{j=0}^{n}a_j(\lambda_j,\omega)\,d\nu(\omega)
\end{equation}
with finite product-seminorm integral.  The corresponding seminorm is the
infimum of that integral over all representations.  After noting the inclusion
of the algebraic tensor product, the authors ask whether the displayed space
is, in general, its completion for these seminorms.

\begin{figure}[ht]
\centering
\includegraphics[width=\textwidth]{figures/open_problem_crop.png}
\caption{Definition 4.1, its seminorm, the algebraic inclusion, and the
completion question on page 22 of the source paper.}
\label{fig:question}
\end{figure}

\section{Proof intuition}

The integral permits a continuum of rank-one factors without requiring the
factor-valued maps to be norm measurable.  Independent signs exploit exactly
that freedom: their correlation is the infinite diagonal matrix.  By contrast,
an algebraic tensor samples to a finite-rank matrix.  The seminorm in the
definition dominates pointwise uniform norm, and the infinite diagonal stays
at least $1/2$ away in uniform norm from every finite-rank matrix.  Thus it
cannot lie in the algebraic completion for the very seminorm used to define
the integral space.

\section{The counterexample}

For $m\geq1$, let
\[
 I_m=[m,m+1),\qquad t_m=m+\tfrac12.
\]
Let $\G$ be the space of all bounded measurable real-valued functions on
$\R$ which vanish outside $[1,\infty)$ and are constant on each $I_m$.
Equip $\G$ with
\[
 \|f\|_{\G,k}:=\|f\|_\infty\qquad(k\in\N).
\]
Thus $\G$ is a Banach function space isometric to $\ell_\infty$, and all its
listed seminorms coincide.

Define
\begin{equation}\label{eq:phi}
 \phi(x,y):=\sum_{m=1}^{\infty}\mathbf 1_{I_m}(x)\mathbf 1_{I_m}(y).
\end{equation}
The sum is pointwise either zero or one.

\begin{theorem}\label{thm:main}
The function $\phi$ belongs to $\G\bi\G$, but for every
$\psi\in\G\otimes\G$ and every pair of seminorm indices $k_0,k_1$,
\[
 \|\phi-\psi\|_{\G\bi\G,k_0,k_1}\geq\frac12.
\]
Consequently, $\G\bi\G$ is not the completion of $\G\otimes\G$ under the
given seminorms.  In particular, the general completion question after
Definition 4.1 of the source paper has a negative answer.
\end{theorem}

\begin{proof}
Let $\Omega=\{-1,1\}^{\N}$ with its product sigma algebra and product
probability measure $\mathbb P$, and write $\varepsilon_m(\omega)=\omega_m$
for its coordinate functions.  Put
\begin{equation}\label{eq:a}
 a(x,\omega):=\sum_{m=1}^{\infty}
 \varepsilon_m(\omega)\mathbf 1_{I_m}(x).
\end{equation}
The map $(x,\omega)\mapsto a(x,\omega)$ is jointly measurable: it is defined
on a countable measurable partition by measurable coordinate functions.  For
every $\omega$, the function $a(\cdot,\omega)$ lies in $\G$ and has all
seminorms equal to one.  Independence and centering of the signs give
\[
 \int_\Omega a(x,\omega)a(y,\omega)\,d\mathbb P(\omega)
 =\sum_{m,l\geq1}\mathbf 1_{I_m}(x)\mathbf 1_{I_l}(y)
   \mathbb E(\varepsilon_m\varepsilon_l)
 =\phi(x,y).
\]
Only one summand in each spatial variable can be nonzero, so the displayed
calculation is pointwise finite.  This is a representation of the form
\eqref{eq:source-rep}, of cost one for every pair $(k_0,k_1)$.  Hence
$\phi\in\G\bi\G$.

We next record that the defining integral seminorm dominates uniform norm.
Indeed, for every representation
\[
 h(x,y)=\int_\Sigma b(x,\sigma)c(y,\sigma)\,d\mu(\sigma)
\]
with factors in $\G$, one has
\[
 |h(x,y)|\leq\int_\Sigma
 \|b(\cdot,\sigma)\|_\infty
 \|c(\cdot,\sigma)\|_\infty\,d|\mu|(\sigma).
\]
Taking first the supremum over $(x,y)$ and then the infimum over
representations yields
\begin{equation}\label{eq:domination}
 \|h\|_\infty\leq\|h\|_{\G\bi\G,k_0,k_1}.
\end{equation}

Now take an algebraic tensor
\[
 \psi(x,y)=\sum_{s=1}^{r}f_s(x)g_s(y)
\]
and sample it at the points $t_m$.  Let
\[
 A_{mn}:=\psi(t_m,t_n),\qquad
 r_m:=(A_{mn})_{n\geq1}\in\ell_\infty.
\]
All rows $r_m$ lie in the finite-dimensional subspace
\[
 F=\operatorname{span}\{(g_s(t_n))_{n\geq1}:1\leq s\leq r\}
 \subseteq\ell_\infty.
\]
Suppose, toward a contradiction, that
\[
 \epsilon:=\sup_{m,n}|A_{mn}-\delta_{mn}|<\frac12.
\]
Then $\|r_m\|_\infty\leq1+\epsilon$ for every $m$.  A bounded subset of the
finite-dimensional normed space $F$ is totally bounded.  Thus there are
distinct $p,q$ such that
\[
 \|r_p-r_q\|_\infty<1-2\epsilon.
\]
Writing $e_m=(\delta_{mn})_{n\geq1}$ for the $m$th identity row gives
\[
 1=\|e_p-e_q\|_\infty
 \leq\|e_p-r_p\|_\infty+\|r_p-r_q\|_\infty
       +\|r_q-e_q\|_\infty
 <2\epsilon+(1-2\epsilon)=1,
\]
a contradiction.  Therefore
\begin{equation}\label{eq:half}
 \sup_{m,n}|\psi(t_m,t_n)-\delta_{mn}|\geq\frac12.
\end{equation}
Because $\phi(t_m,t_n)=\delta_{mn}$, equations
\eqref{eq:domination} and \eqref{eq:half} prove
\[
 \|\phi-\psi\|_{\G\bi\G,k_0,k_1}
 \geq\|\phi-\psi\|_\infty\geq\frac12.
\]
Thus no net or sequence of algebraic tensors can converge to $\phi$ for the
given seminorms.
\end{proof}

\section{Verification, scope, and novelty}

The construction checks every condition in Definition 4.1 directly: the
parameter measure is finite, the factor is jointly measurable, every section
belongs to $\G$, the norm map is the constant one, and the product-norm
integral equals one.  The separation uses the exact seminorm from that
definition, not the larger Banach projective norm.  Since that seminorm
dominates uniform norm here, it is definite and no quotient-completion issue
arises.

This fully answers the words ``in general'' in the negative.  It does not
decide density for the special spaces $S^\beta(\R)$, $T^\beta(\R)$, or the
spectral $L_\infty^\beta$ spaces used elsewhere in the source.

Before promotion, the run indexes and ledgers were searched using the source
arXiv id and the core phrases ``integral projective tensor'', ``algebraic
tensor completion'', ``Rademacher'', and ``diagonal''.  A bounded web/arXiv
search used the exact question sentence, the source title and id, and close
integral-projective-tensor phrases.  No separate answer was found.  Moreover,
Hekkelman's 2026 thesis repeats the question verbatim in Remark 3.2.2
\cite{HekkelmanThesis}.  This is strong recent status evidence, not an
exhaustive novelty guarantee.

Human review is recommended.  The highest-value checks are the interpretation
of completion in the source sentence and the total-boundedness separation of
the sampled rows.

\begin{thebibliography}{9}
\bibitem{HMN}
E.-M. Hekkelman, E. McDonald, and T. D. H. van Nuland,
\emph{Multiple operator integrals, pseudodifferential calculus, and asymptotic
expansions}, arXiv:2404.16338, 2024.
\url{https://arxiv.org/abs/2404.16338}.

\bibitem{HekkelmanThesis}
E.-M. Hekkelman,
\emph{Trace Formulas}, doctoral thesis, 2026, Chapter 3, Remark 3.2.2.
\url{https://www.emhekkelman.nl/uploads/phdthesis.pdf}.
\end{thebibliography}

\end{document}

